\section{Poisson disintegration for patch factorization}
\label{app:patch-factorization-proof}

This section gives the measure-theoretic justification of
Theorem~\ref{thm:patch-factorization}. The main text records the product
conditioning mechanism; here we justify the conditional law at fixed successful
interaction times.

\begin{proof}[Detailed proof of Theorem~\ref{thm:patch-factorization}]
Write~$G_T=(Y_0,\mathcal I_T)$, so~$\mathcal G_T=\sigma(G_T)$. By
Lemma~\ref{lem:dual-local-finiteness}, its positive-time records under~$\pP_A$
almost surely form a finite chronological list
\begin{equation*}
g=\bb{(i_k,t_k,S_k)}_{k=1}^n,
\qquad
0<t_1<\cdots<t_n\leq T,
\end{equation*}
where~$\varnothing\neq S_k\subseteq N(i_k)$ and
$i_k\in A\cup\bigcup_{\ell<k}S_\ell$. The list space and the finite
marked-configuration spaces are standard Borel.

On the component with fixed discrete records~$((i_k,S_k))_{k=1}^n$, set
\begin{equation*}
m_T(dg)
=
\prod_{k=1}^n
\bb{\delta_{i_k}(S_k)+\beta_{i_k}(S_k)}
\,dt_1\cdots dt_n,
\end{equation*}
and sum these measures over~$n\geq0$ and the discrete records, omitting
zero-rate components. For nonnegative measurable~$h$ and~$f_P$, the
multivariate Mecke formula~\cite[Theorem~4.4]{LastPenrose2017}, applied to
ordered tuples of nonempty-target marks, gives
\begin{equation}
\label{eq:appendix-patch-factorization-mecke}
\E_A\Cb{
h(G_T)
\prod_{P\in\mathcal P_T}f_P\bb{\Sigma_P}
}
=
\int h(g)
\prod_{P\in\mathcal P_T(g)}
\E_P\Cb{
f_P\bb{\Sigma_P}\ind_{\operatorname{Con}(P)}
}
\,m_T(dg).
\end{equation}
Indeed, a mark with record~$(i,t,S)$ has combined intensity
$\bb{\delta_i(S)+\beta_i(S)}\,dt$ and, conditional on this record, its kind has
the law in~\eqref{eq:patch-initial-kind-law}. The remaining marks in the
disjoint patch intervals are independent with laws~$\pP_P$.

It remains to identify the consistency indicators in
\eqref{eq:appendix-patch-factorization-mecke}. Suppose first that every patch is
consistent and read the records chronologically. The initial incoming patches
give the active set~$A$. Between record times, empty-target marks change the
local and global active indicators in the same way, while consistency makes
every unrecorded nonempty-target mark unsuccessful. At a recorded interaction,
consistency makes its source active; its kind determines the new state of the
source, and its targets start the incoming patches. Hence the prescribed
records are exactly the successful-interaction skeleton. Conversely, if the
prescribed records are the complete successful-interaction skeleton, the
one-site processes~$X^P$ are the restrictions of the global active process and
every patch is consistent.

Taking~$f_P=1$ in~\eqref{eq:appendix-patch-factorization-mecke} gives
\begin{equation*}
\pP_A\bb{G_T\in dg}
=
\prod_{P\in\mathcal P_T(g)}
\pP_P\bb{\operatorname{Con}(P)}
\,m_T(dg).
\end{equation*}
Taking the Radon--Nikodym derivative with respect to this marginal yields, for
$\pP_A\circ G_T^{-1}$-almost every~$g$,
\begin{align*}
\E_A\Cb{
\prod_{P\in\mathcal P_T}f_P\bb{\Sigma_P}
\,\middle|\,G_T=g
}
&=
\prod_{P\in\mathcal P_T(g)}
\frac{
\E_P\Cb{f_P\bb{\Sigma_P}\ind_{\operatorname{Con}(P)}}
}{
\pP_P\bb{\operatorname{Con}(P)}
}
\\
&=
\prod_{P\in\mathcal P_T(g)}
\E_P^{\mathrm{con}}\Cb{f_P\bb{\Sigma_P}}.
\end{align*}
This proves~\eqref{eq:patch-factorization} for nonnegative factors and
identifies the conditional law. The integrable signed case follows from this
conditional-law statement.
\end{proof}
